\chapter{Patterns}
\label{ch:patterns}

A pattern is a special language construct that describes a structure of data.

A pattern is used to perform the "deconstruction" of a value into the structure
that the pattern describes.

\begin{minted}{\giovilexer}
getPoint() => tuple<Int, Int> {
	=> <69, 420>
}
<x, y> <- getPoint()
"The point is at $x and $y" -> _out
\end{minted}

Patterns can also be used as iterators in a for loop; at each iteration the
value will be matched to the pattern.

\begin{minted}{\giovilexer}
const CITIES <- [
	"Rome"    => <41.8941047, 12.5151106>
	"NewYork" => <40.7157121, -74.0013879>
	"Tokyo"   => <35.7104635, 139.7300523>
]
for city => <x, y> in CITIES do {
	"$city is located at $x and $y" -> _out
}
\end{minted}

\paragraph{Patterns compatibility} A value can be deconstructed into a pattern
only when it is "equal or narrower" than the pattern; the rules are similar to
types widening rules.

\todo{
	mmhhh... questo non è esattamente corretto. $\!x:\Int\rightarrow p_x:\Real$ è
	compatibile ma,
	$\!\tuple{x,y}\rightarrow\tuple{p_x,p_y,p_z}$ non lo è
}

\begin{minted}{\giovilexer}
getPoint() => tuple<Int, Int> { => <69, 420> }
getPoint() -> <x Real, y Real>
"$x and $y are Real" -> _out
\end{minted}

\clearpage

\section{Patterns kinds}

There are five kinds of patterns:
\begin{itemize}
	\item Terminal patterns:
		\begin{itemize}
			\item \textbf{Placeholder pattern}: \code{\_} (underscore)
				\todo{rinominare. Sinkhole?}
			\item \textbf{Unary pattern}: \code{someIdentifier}
		\end{itemize}
	\item \textbf{Product pattern}:
		\begin{itemize}
			\item \textbf{Positional}: \code{<a, b, c>}
			\item \textbf{Nominal}: \code{<x:a, y:b, z:c>}
		\end{itemize}
	\item Iterator patterns:
		\begin{itemize}
			\item \textbf{List pattern}: \code{[someIdentifier]}
			\item \textbf{Map pattern}: \code{[key => value]}
		\end{itemize}
\end{itemize}

\section{Nesting patterns}

Non-terminal patterns can have other patterns nested.\\*
An extreme example: a value of type \code{[Text=>[tuple<Int,Int,Int>]]} (map
with $\Text$ keys and a list of tuples as value) could be iterated with a
pattern \code{[colorName=>[<r,\_,b>]]} is a map pattern with a unary pattern as
key and a list pattern as value which itself has a positional product pattern
nested, where the second field pattern is a placeholder.

\nb{
	Nesting iterator patterns effectively means nesting iterations.\\*
	For example \code{for [[[i]]] in matrix3D do} would be equivalent of
	writing three nested for loops
}

\subsection{Limitations}

\begin{itemize}
	\item Product patterns can only have terminal patterns (placeholder, unary)
		or other product patterns nested.
	\item In map patterns, the pattern for the key can only be a terminal or a
		product pattern; the pattern for the value can be of any kind (other
		iterator patterns as well).
\end{itemize}

\clearpage

\section{Terminal patterns}

A terminal pattern is called as such because it instructs not to deconstruct
the value further.

\subsection{Placeholder pattern}

The simplest kind of pattern is the placeholder pattern (denoted by an
underscore \code{\_}).\\*
It basically acts as a sinkhole; redirecting to it does nothing. It can be
useful to explicitly discard a function return value, or when nested in other
patterns.

\begin{minted}{\giovilexer}
stack.pop() -> _ # pop a stack but ignore the value

getPoint() -> <_, y>
"I only care about $y" -> _out
\end{minted}

\subsection{Unary pattern}
A unary pattern is a terminal pattern that simply states a single "unit" to
which the value is assigned to as a whole, without deconstruction.\\*
The pattern is denoted by an identifier (the name) and an optional type to
which the deconstructed value is widened to.\footnotemark[1]

\footnotetext[1]{A variable declaration is secretly a unary pattern!}

\clearpage

\section{Product patterns}

A product pattern describes the data structure of a product type, hence it is
meant to be used with product values (tuples, records, structs).

Like mentioned above, the patterns follow the same rules of types
compatibility; in order to be able to deconstruct all product values, you are
given the choice of two deconstruction approaches: positional and nominal.

\subsection{Positional}

With a positional product pattern, the fields of a product value are
deconstructed in the same order they're defined in the pattern.

The positional product pattern can be used to deconstruct all product types.

\begin{minted}{\giovilexer}
<x, y> <- getPoint()
"The point is at $x and $y" -> _out
\end{minted}

\subsection{Nominal}

With the nominal approach you're deconstructing the value based on the fields'
names (the order won't matter), therefore the pattern labels must be the same
of the fields' names.

Obviously, nominal product patterns can only be used with product types with
named fields (i.e. records, structs).

\begin{minted}{\giovilexer}
<y:someY, x:someX> <- getPoint()
"The point is at $someX and $someY" -> _out
\end{minted}

\subsubsection{Implicitly named product pattern}

It is often needed to deconstruct a nominal product value into a pattern whose
variables share the same names.\\*
A shorthand syntax is available for such scenario; however, only unary patterns
are allowed. In this case, the identifier of the unary pattern will be used for
the field lookup of the product value

\begin{minted}{\giovilexer}
<_:y, _:x> <- getPoint()
"The point is at $x and $y" -> _out
\end{minted}

\clearpage

\section{Iterator patterns}

Iterator patters can only be used in a for loop statement (hence the name). For
each iteration of the loop, the iterated value is deconstructed into the
specified pattern.

\subsection{List pattern}

The list pattern is used to iterate over a list; it's basically a unary
pattern, but allows nested patterns.

\begin{minted}{\giovilexer}
# iterate elements of a list
for [i] in [0 1 2 3] do somethingWith(i)

# iterate elements of a nested list
identityMatrix [[Real]] <- [
	[1 0 0]
	[0 1 0]
	[0 0 1]
]
for [[i]] in identityMatrix do somethingWith(i)
\end{minted}

Since it is very common to iterate the elements of a list, the brackets can be
omitted for non-nested patterns

\begin{minted}{\giovilexer}
for i in [0 1 2 3] do somethingWith(i)
\end{minted}

\subsection{Map pattern}

The map pattern is used to iterate over maps; it's a pair of two patterns, one
for the key and one for the value respectively.

\begin{minted}{\giovilexer}
# iterate over a map
for [key => value] in ["Simone"=>69 "Luca"=>420 "Paolo"=>10] do {
	somethingWith(key, value)
}

# iterate over a nested map
graph [Text => [Text => Int]] <- [
	"a" => ["b"=>5 "e"=>3]
	"b" => ["a"=>5 "c"=>1 "d"=>9 "f"=>2]
	"c" => ["b"=>1 "e"=>6]
	"d" => ["b"=>9 "e"=>7]
	"e" => ["a"=>3 "c"=>6 "d"=>7 "f"=>8]
	"f" => ["b"=>2 "e"=>8]
]

for nodeA => [nodeB => edgeCost] in graph do {
	"$nodeA ---$edgeCost---> $nodeB" -> _out
}
\end{minted}

Like for the list pattern, the brackets can be omitted for non-nested patterns

\begin{minted}{\giovilexer}
for key => value in ["Simone"=>69 "Luca"=>420 "Paolo"=>10] do {
	somethingWith(key, value)
}
\end{minted}

